% ============================================================
%  R. Nielsen, "Om Theologiens Naturbegreb, med særligt Hensyn til
%  Malebranche: De la recherche de la vérité"
%  Kjøbenhavn: Schultziske Officin, 1855
%  TRANSLATION (English) — Hans Halvorson
% ============================================================
\documentclass[12pt,a4paper]{article}
\usepackage[utf8]{inputenc}
\usepackage[T1]{fontenc}
\usepackage{libertinus}
\usepackage{libertinust1math}
\usepackage{textalpha}   % polytonic Greek (p.34: Aristotle quote, ὁμοούσιος, θεία φύσις)
\usepackage{enumitem}    % in case any enumerated lists are needed
\usepackage[english]{babel}
\usepackage[protrusion=true,expansion=true]{microtype}
\usepackage{setspace}
\usepackage[margin=1.2in]{geometry}
\usepackage{fancyhdr}
\usepackage{hyperref}

\hypersetup{
  pdftitle={On Theology's Concept of Nature},
  pdfauthor={R. Nielsen},
  hidelinks
}

\pagestyle{fancy}
\fancyhf{}
\rhead{Nielsen, \emph{On Theology's Concept of Nature} (1855)}
\cfoot{\thepage}

\setstretch{1.3}

\title{\textbf{On Theology's Concept of Nature}\\[4pt]
       \large with Special Reference to\\[4pt]
       \large Malebranche: \emph{De la recherche de la vérité}}
\author{R.\ Nielsen\\[4pt]
  \small Translated from the Danish by Hans Halvorson}
\date{Copenhagen: Schultz, 1855\\[4pt]
  \small Invitation to the University of Copenhagen's Annual Festival\\
  \small in Memory of the Reformation of the Church}

\begin{document}
\maketitle
\thispagestyle{fancy}

\bigskip

% === TRANSLATION BEGINS HERE ===
% Source of truth: transcription.tex (Danish, COMPLETE). Translate FROM it, not
% from the scan. Work the markers below in reading order, ~10 printed pages per
% batch; replace each marker with the English. See TRANSLATION-HANDOFF.md for the
% terminology glossary and the rules on quotes / French footnotes / Greek / Latin.
% The essay proper is printed pp. 3--37; pp. 38--47 are the appended vitae +
% festival invitation (translate only if in scope — see the hand-off doc).

% ---- The treatise (printed pp. 3--37) ----
% p. 3
\noindent
Between theology and natural science there is no understanding; theology's
concept of nature is not natural, and natural science's is not theological.

Even by merely confining oneself to a factual juxtaposition of those natural
events in whose interpretation natural science most decidedly contradicts
theology, it would be easy to make the disagreement between the two sciences
conspicuous. Touch whichever dogma, whichever theological presupposition one
will; insofar as there can in any way be a question of a real, corresponding
natural relation, natural science, by virtue of the investigation of the facts,
must reject the theological explanation.

Is it the creation in six days: natural science proves that for the formation of
the earth alone a period of more than $6 \times 6$ million years must have
elapsed.

Is it the dogma of Paradise and the Fall: natural science proves that there was
corruptibility, pain, and death in the world long before the first human beings
came to be.

Is the talk of the Flood and Noah's Ark: natural science proves that ``the
Flood'' is a natural local catastrophe, and that the animals could not possibly
have been saved in Noah's Ark, since the beasts of prey would in a few days have
devoured the other animals, and thereafter have died of hunger.

Thus one could go on, and still not reach the goal. Even if the juxtaposition
were carried through ever so completely, the grounds and counter-grounds from
both sides sought out and weighed ever so carefully, it would nevertheless along
this path hardly be possible to arrive at any decisive result.

% p. 4
To be sure, natural science, with respect to the main points of the matter, is
on the whole in agreement with itself, just as theology too must on the whole be
said to follow a direction answering to its own peculiar endeavours; but one may
not on that account suppose that either of the two sciences would at a given time
be wholly and entirely in agreement with itself in the interpretation of
particulars.

Natural science is, as they say, in becoming; this becoming --- under which, as
regards particulars, it continually corrects itself, secures its method,
purifies and broadens its foundation --- is precisely characteristic of its
luxuriant growth, its restless, imposing progress.

Theology, then, is in its own way also in becoming --- a becoming, namely, under
which its ever-increasing lack of principle and the inward discord of its members
point more or less toward the weakness of old age and impending
dissolution\footnote{Lest it should seem to some reader less acquainted with the
discussions belonging here as though this assertion were thrown out without any
foundation, I must for the present refer to my work \emph{Evangelitroen og
Theologien}. I may do so with all the more right, since no theologian yet, so far
as is known to me, has found himself moved to show the groundlessness of the view
set forth in the said work.}.

If both sciences are now thus each in its own becoming, then it would surely help
but little for anyone to begin afresh either to weigh the grounds for and against
``the Flood,'' or to debate the problem of ``Noah's Ark,'' or, in other words, to
set natural science's conception of particular phenomena against theology's
particular dogmas.

If an understanding between theology and natural science is to be thought
possible, it must presumably begin with the two sciences coming to agreement in
the determination of that concept which, for the scientific reality of both, is a
decisive fundamental concept, namely the concept: \emph{nature}.

By such a determination of the concept one cannot here first of all think of some
definition in general; for with formal definitions little
% p. 5
is accomplished in natural science. The task is to indicate the point of view
under which the essence of nature can be seen in the purest, truest scientific
light.

Relying on the testimony of the senses and the validity of thought, natural
science begins by taking existing things as they are given, then advances by
exact examination from the known to the as-yet-unknown, and everywhere confines
itself to seeing nature in nature's light.

Theology, on the other hand, chooses for itself a higher point of departure.
Beginning with the Creator, it passes through creation over to the creature, and
so comes, by its own opinion, to behold nature as God's work in God's light.

If it now really is the case that there truly exists a science able to present
nature in God's light, and if this science is theology: then indeed the
theological concept of nature, whatever may be said for or against the
interpretation of any single phenomenon, is undeniably the highest norm by which
empirical science, in its judgment of the essence of nature, has to direct
itself.

But should it, on the other hand, prove scientifically demonstrable that there
neither is nor can be a science able in any way, in any degree, to elucidate the
natural by means of the supernatural: might we not then conclude from this that
the theological doctrine of creation, with its corresponding explanation of
nature, had best be assigned a place of its own outside of science?

As a representative of the theological conception of the essence of nature, we
choose, for the clarification of the matter, a keen thinker, the author of the
work \emph{De la Recherche de la Vérité}\footnote{Paris MDCLXXVIII.}, the noble
Malebranche.

That Malebranche is a Catholic, and that his Cartesian natural philosophy has
long since been antiquated, is of no significance in this connection.

Here, as regards the dogmas, there is only one dogma on which it essen%
% p. 6
tially depends; it is the dogma of creation. Here, as regards natural science,
there is only one task on which attention is fixed; it is the task of recognizing
nature as God's work, of beholding all things in God.

Inwardly convinced that this great task can be solved, M.\ is so far from
withdrawing into the mysterious, and concealing poverty of thought in
aesthetically nebulous speculations, that he rather calls for beginning with the
particular, the distinct, the easily comprehensible, and in many places speaks a
language such that one would believe it was not one of the Fathers of the Oratory
speaking of God's omnipotence, but one of today's naturalists dismissing
philosophical mysticism\footnote{Jfr. t.~Ex. De l. Recherch. S.~394. Combien de
gens rejettetent la Philosophie de M. Descartes par cette plaisante raison que
les principes en sont trop simples et trop faciles. Il n'y a point de termes
obscurs et mysterieux dans cette Philosophie: des femmes et des personnes qui ne
sçavent ni grec ni latin, sont capables de l'apprendre: il faut donc que soit peu
de chose, et il n'est pas juste que de grands genies s'y appliquent. Ils
s'imaginent que des principes si clairs et si simples ne sont pas assez féconds,
pour expliquer les effets de la nature qu'ils supposent obscure et embarassée.
\dots\ Ils aiment donc mieux expliquer des effets, dont ils ne comprennent point
la cause, par des principes qu'ils ne conçoivent point, et qu'il est absolument
impossible de concevoir, que par des principes simples et intelligibles
tout-ensemble. Car ces Philosophes expliquent des choses obscures par des
principes qui ne sont pas seulement obscurs, mais entièrement
incompréhensibles.}.

\begin{center}---\end{center}

``It is'' --- so thinks Malebranche --- ``error that is the cause of man's misery;
it is therefore necessary that it be combated with all one's powers. One must not
imagine that it is so troublesome to seek the truth. One need only open one's
eyes, be attentive, and follow a few rules which we shall give in what follows.
For it is, after all, sharp thinking that brings light and discovers the truth.''

``Although our senses are corrupted by sin, it is nevertheless not really the
senses, but our freedom, that is the cause of our errors.''

% p. 7
``For since we are composed of mind and body (Etant composez d'un esprit et d'un
corps), we have two kinds of goods to seek, the mind's and the body's. We have
then also two ways in which we can know whether something is good or evil. We can
know a good either by a clear insight or by an obscure sensation (un sentiment
confus). I recognize by reason that justice is lovable; I also know by taste that
this or that fruit is good. The beauty of justice cannot be sensed (ne se sent
pas), the goodness of the fruit cannot be thought (ne se connoit pas). The goods
of the body therefore do not deserve that a mind, which God has created only for
himself, should scrutinize them; it is therefore necessary that the mind, without
further examination (examen), be able to know such goods by the short,
indubitable test of sensation\footnote{S.~18. Les pierres ne sont pas propres à
la nourriture, la preuve en est convaincante, et le seul goût en a fait tomber
d'accord tous les hommes.}.''

``If the mind saw in bodies only what is really in them, then it could not love
them, and could not make use of them without much reluctance; it is therefore
necessary that they seem agreeable, in that they occasion feelings that do not
pertain to their essence.''

``But it is not so with God, who alone is the mind's true good; for him it is not
enough to be loved with the blind love of instinct; he wills to be loved with an
enlightened, clear, and resolved love (d'un amour éclairé et d'un amour
déchoix).''

``To understand this rightly, one must know that there is a very considerable
difference between the way in which the Author of Nature (l'Auteur de la Nature)
moves matter, and the way in which he without ceasing moves the mind toward the
good in general (vers le bien en general). For matter is without any activity,
without power to arrest its motion or determine its direction.''

``It is otherwise with the will; of it one can nevertheless in a certain sense
(en un sens) say that it is active, and that it has in itself the power to take up
the impression (impression) God gives
% p. 8
it\footnote{S.~5. car quoi qu' elle ne puisse pas arrêter cette impression, elle
peut en un sens la détourner du côté qu'il lui plaît\dots} in various ways, and
thereby to occasion all the disorder (déreglement) that shows itself in the
inclinations, and all the miseries that are consequences of sin.''

``Even if the will is not free in the sense that it should have the power to will
and not to will, it is nevertheless able to drive us forward toward the good in
general\footnote{S.~5. \dots\ par ce mot de \emph{Volonté}, je prétens ici
désigner l'impression ou le mouvement naturel, qui nous porte vers le bien
indéterminé et en général: et par celui de \emph{Liberté}, je n'entens autre
chose que la force qu' a l'esprit de détourner cette impression vers les objets,
qui nous plaisent.}.''

``Since the will always directs itself according to the representation of the
good, and we cannot possibly will that which at the same moment we must regard as
an essential evil, there must be the most intimate connection between cognition
and will, between the true and the good\footnote{S.~6. C'est à-peu-près la même
chose de la connoissance de la vérité que de l'amour du bien.}.''

``The senses and passions (les sens et les passions) do not have their origin
from sin, but the power with which they tyrannize sinners does. This power is not
so much a disturbance on the side of the senses as a confusion in man's mind and
will.''

``Adam, in his state of innocence, had the same senses as we, but was the
unconditional master of the impressions of joy and pain that they awakened in his
body. Only when he voluntarily turned away from the countenance of God, as it
were to sate his mind with the beauty and anticipated sweetness of the forbidden
fruit, did his senses and passions rebel against him, and make him, like us, a
slave of all sensuous things.''

``Our senses are not so corrupted as one imagines; but it is the innermost of our
soul, it is our freedom, that is corrupted. It is
% p. 9
not our senses that deceive us; it is our will that deceives us by its rash
judgments.''

``When one, for example, sees light, then it is quite certain that one sees
light; when one feels warmth, one does not deceive oneself in believing that one
feels it --- be it before or after the Fall.''

``But one deceives oneself when one assumes (quand on juge) that the warmth one
feels is outside the soul that feels it. It is a rash judgment, resting on the
misuse of freedom. Therefore we ought carefully to observe the rule never to
judge by the senses what things are in themselves; for the senses are in reality
not given that we should know the truth of things in themselves, but only to
preserve our body (pour la conservation de notre corps).''

``The ground of error also does not lie in the limitation of our understanding;
for the understanding is indeed limited, but it is nevertheless the will that
guides our power of judgment\footnote{S.~6 \dots\ nous sommes aussi libres dans
nos faux jugemens, que dans nos amours déreglez\dots}. It therefore becomes a
task of the will to work its way out of the illusion and seek the truth. But to
seek the truth is in all things to seek God.''

``Truth is known by its clarity, by its self-evidence\footnote{S.~9. Car la
vérité ne se trouve presque jamais qu' avec l'évidence, et l'évidence ne consiste
que dans la vûe claire et distincte de toutes les parties et de tous les rapports
de l'objet, qui sont nécessaires pour porter un jugement assûré.}; it must
therefore be impressed upon us that we never grant anything before it is so
evident to us that we are as it were compelled by the inner, earnest correction
of our reason (jusqu'à ce que nous y soyons comme forcez par des reproches
intérieures de notre raison\footnote{S.~1--2. L'exactitude de l'esprit n'a presque
rien de penible: ce n'est point une servitude comme l'imagination la représente;
et si nous y trouvons d'abord quelque difficulté, nous en recevons bientôt des
satisfactions qui nous récompensent bien de nos peines; car enfin il n'y a
qu'elle qui produise la lumiére, et qui découvre la vérité.}).''

% p. 10
After thus having pointed out the cause of error, and having besides (in the
first and second Books) treated in detail the errors occasioned by the senses and
the imagination, the author passes over (in the third Book) to treat of the
essence of the mind: \emph{De l'entendement ou de l'esprit pur}.

The general relation between mind and matter has already, in the first Book, been
conceived in the following way:

``Just as matter, the extended, has a twofold capacity --- a capacity to receive
different form, and a capacity to be moved --- so too the human mind has two
faculties, \emph{the understanding} (l'entendement), by which we apprehend
different ideas, that is: bring different things to consciousness, and \emph{the
will} (la volonté), the capacity for manifold inclinations, or the willing of
different things.''

This juxtaposition of the material and the mental, whereby apparent points of
resemblance are brought out, must not, however, be understood as though a natural
kinship between mind and matter were here presupposed. On the contrary! It is
impressed throughout the whole investigation, from first to last, that the two
substances are as heterogeneous as possible.

Although different faculties are ascribed to them, they nevertheless have no
active faculties\footnote{S.~4. C'est pour nous accommoder à la maniére ordinaire
de parler, que nous dirons dans la suite, que les sens sentent\dots}. The
concept ``faculty'' must here, for the mind as well as for matter, be taken in so
passive a sense that it does not at all express a capacity for reciprocal
activity, but only a possibility of being acted upon by God\footnote{S.~5. De même
que l'Auteur de la Nature est la cause universelle de tous les mouvemens, qui se
trouvent dans la matiére; c'est aussi lui qui est la cause générale de toutes les
inclinations naturelles qui se trouvent dans les esprits\dots}.

``The essence of the mind,'' it now says (third Book), ``consists in thinking
alone, just as that of matter in extension alone. It is in accordance with the
different modifications of thinking that the mind is now willing, now
representing, or has several other peculiar forms.''

% p. 11
``It likewise depends on different modifications, namely of extension, that matter
is now water, now wood, now fire, or has an infinity of other such
forms\footnote{S.~171. Je ne croi pas aussi qu'il soit possible de concevoir un
esprit qui ne pense point, quoi qu'il soit fort facile d'en concevoir un qui ne
sente point, qui n'imagine point, et même qui ne veüille point: de même qu'il
n'est pas possible de concevoir une matiére, qui ne soit pas étendue quoi qu'il
soit assez facile d'en concevoir une, qui ne soit ni terre, ni métal ni quarrée
ni ronde et qui même ne soit point en mouvement. Il faut conclure de là, que
comme il se peut faire qu'il y ait de la matiére, qui ne soit ni terre ni metal,
ni quarrée ni ronde ni même en mouvement: il se peut faire aussi qu'un esprit ne
sente ni chaud ni froid, ni joie, ni tristesse, n'imagine rien, et même ne veuille
rien; de sorte que toutes ces modifications ne lui sont point essentielles.
\emph{La pensée toute seule est donc l'essence de l'esprit, ainsi que l'étendue
toute seule est l'essence de la matiére.}}.''

Since now the essence of the mind is placed in thought, it might indeed seem as
though on the whole too strong an emphasis were here laid on the will; but M.\
expressly remarks that the faculty of willing is nevertheless inseparable from the
mind, although it does not constitute its essence, just as the faculty of being
moved is inseparable from matter, although it likewise does not constitute
matter's essence.

``Even if thinking is the essence of the mind, we may nevertheless by no means
conclude from this that the thinking mind should be able of itself to know all the
modifications whose possibility it contains\footnote{S.~173. Un simple morceau de
cire est donc capable d'un nombre infini, ou plûtôt d'un nombre infiniment infini
de différentes modifications, que nul esprit ne peut comprendre. Quelle raison
donc de s'imaginer que l'âme qui est beaucoup plus noble que le corps, ne soit
capable que des seules modifications qu'elle a déja reçuës.}.''

``When our soul here below receives only very few modifications, it is because it
is bound to a body and dependent on it.''

``The matter of which our body is composed is able
% p. 12
to receive only very few modifications, as long as we live. This matter cannot
dissolve into earth, into vapour, until after our death. During our lifetime it
cannot become air, fire, diamond, metal; it cannot become round, square,
triangular: it must be flesh, and have a human figure, in order that our soul may
be united with it.''

``It is the same with the soul. It is necessary that it have sensations of warmth,
of cold, of colour, of light, of tones, of smell, of taste, and many other
modifications. These sensations teach it precisely to preserve its machine
(l'appliquent à la conservation de sa machine\footnote{Elles l'agitent, et
l'effraient dés que le moindre ressort se débande, ou se rompt, et ainsi il faut
que l'âme y soit sujette, tant que son corps sera sujet à la corruption. Mais
lorsqu'il sera revêtu de l'immortalité, et que nous ne craindrons plus la
dissolution de ses parties, il est raisonable de croire qu'elle ne sera plus
touchée de ces sensations incommodes que nous sentons malgré nous; mais d'une
infinité d'autres toutes différentes, dont nous n'avons maintenant aucune
idée.}).''

``Let one beware of dispersing the mind by wanting to explain things that are far
too intricate. For the human mind is subject to errors, not only because it is not
infinite, but also because it does not have enough perseverance to dwell on each
single object so long, until it has been in every respect completely and exactly
investigated.''

``The objects with which the mind occupies itself can be divided into two classes:
they are either in the soul or outside it.''

``Those that are in the soul are its own thoughts, by which are in general
understood all the things (toutes les choses!) that cannot be in the soul without
its becoming aware of them, such as: its own sensations, its imaginations, pure
concepts and representations, even its passions (ses passions mêmes) and natural
inclinations.''

As regards these objects, the soul can indeed know them without needing an
intermediary; for they are nothing other than the soul itself, modi%
% p. 13
fied in this or that way (d'une telle ou telle façon); but \emph{how does the soul
come to the knowledge of the objects that are outside the soul?}

That is a chief question.

``We see the sun, the stars, an infinite multitude of objects outside us; yet it
is not probable that the soul goes out of the body, that it, so to speak, takes a
stroll through the heavens in order there to behold all these objects. Nor, then,
does it see them by themselves. The mind's immediate object, when it sees the sun,
is not the sun itself, but a something that is intimately united with our soul,
and it is this that I call the idea'' (S.~188. \emph{De la nature des idées}).

The characteristic feature of M.'s doctrine of ideas is now a direct consequence
of the peculiar way in which he, in relation to the philosophical standpoint of
the time, explains to himself the knowledge of external objects.

``\emph{Material objects cannot possibly be apprehended by themselves; they are
apprehended only by the ideas corresponding to them in the soul: whence then come
these ideas?}\footnote{S.~190. Nous assurons donc qu'il est absolument nécessaire,
que les idées, que nous avons des corps, et de tous les autres objets, que nous
n'appercevons point par eux mêmes, viennent de ces mêmes corps ou de ces objets:
ou bien que nôtre âme ait la puissance de produire ces idées: ou que Dieu les ait
produites avec elle en la créant, ou qu'il les produise toutes les fois qu'on
pense à quelque objet: ou que l'âme ait en elle-même toutes les perfections
qu'elle voit dans ces corps: ou enfin qu'elle soit unie avec un être tout parfait,
et qui renferme généralement toutes les perfections des êtres créez.}''

Toward an answer to this question M.\ makes his way through a critique of the
theories diverging from his own view.

``The Peripatetics' opinion, that external objects should send us their images
(des espèces, qui leur ressemblent), and that these should then, through the outer
senses, be brought to the common sense (au sens commun),
% p. 14
is unreasonable; these impression-images (especes impresses) would then themselves
have to be little bodies\footnote{S.~191. Tous les objets, comme le Soleil, les
Etoiles, et tous ceux qui sont proche de nos yeux, ne peuvent pas envoyer des
especes, qui soient d'autre nature qu'eux \dots\ c'est pourquoi les Philosophes
disent ordinairement, que ces especes sont grossieres et matérielles, à la
différence des especes \emph{expresses}: (especes matérielles et sensibles, qui
sont rendues intelligibles par l'intellect agent ou agissant, et sont propres pour
être reçues dans l'intellect patient S.~190) qui sont spiritualisées.}, and these
bodies, with which space would be filled, would then, owing to the impenetrability
of matter, have to crush and collide with one another, some going to one side,
others to another, and thus could not make the objects visible.''

The opinion of those who assume that the ideas of objects are created with us, so
that God does not need at each moment to awaken them in the soul when we have need
of them, is refuted by the following consideration:

``We have, even as regards simple figures, an infinite multitude of ideas.
Ellipses, triangles, polygons, and the like can, with respect to proportions of
magnitude, be imagined in innumerable ways\footnote{S.~196. L'esprit voit donc
toutes ces choses: il en a des idées. Il est sûr que ces idées ne lui manqueront
jamais, quand il employeroit des siécles infinis à la considération même d'une
seule figure.}. It is then highly improbable that God should have created such a
swarm of ideas in and with the human mind. But suppose that we did have such a
store of co-created ideas (un magazin de toutes les idées), it still remains
wholly inexplicable how we come to apply the right idea when the corresponding
object presents itself. We see, for example, the sun. The image that the sun
imprints in the brain has, as already shown, no resemblance to the idea we form of
the sun; nor does the soul notice the least thing of the impression the sun makes
either on the ground of the eye or in the brain: it is therefore impossible to
explain how, among the innumerable ideas, we come precisely to choose the idea of
the sun when we see the sun.''
% p. 15
``Those who maintain that our souls themselves have the faculty of being able to
form ideas corresponding to the objects about which we would think, although the
impressions coming from these objects are not images resembling the objects, have
either far too poor conceptions of ideas or far too lofty thoughts of the human
mind\footnote{S.~193. En effet le monde intelligible doit être plus parfait que le
monde matériel et terrestre\dots\ Ainsi, quand on assure, que les hommes ont la
puissance de se former les idées telles qu'il leur plaît, on se met fort en danger
d'assurer que les hommes ont la puissance de faire des êtres plus nobles et plus
parfaits que le monde que Dieu a créé. On ne fait pas cependant réflexion à cela,
parce qu'on s'imagine, qu'une idée n'est rien, à cause qu'elle ne se fait point
sentir \dots}.''

``Such an imagining of the soul's perfection is presumptuous. God alone can in
himself behold and know all things, since he is the Creator of all things. `Say
not that you have the light in yourselves,' says St.\ Augustine, `for God alone
has the light in himself.'\,''

``There then remains only one view as the sole rational one (qui paroit seule
conforme à la raison), namely, that minds have all their thoughts from God.''

``To grasp this rightly, we must first recall that God necessarily must have the
ideas of all created things in himself, since otherwise he could not produce them;
next, that God, by virtue of his omnipresence, is so intimately united with our
souls that one could even say that he is the place of minds ((qu'il est le lieu des
esprits), just as one says of space that it is the place of bodies.''

``Under these presuppositions it is therefore certain that the mind can see what
it is in God that presents created things, since it is so spiritual, so
intelligible, so present to the mind. The mind can then see in God God's work, so
truly as God wills to reveal it to it.''
% p. 16
Thus M.\ comes to set up the principle so indispensable for theology, namely, that
we behold all things in God:

\begin{center}
\textbf{Que nous voyons toutes choses en Dieu.}
\end{center}

This fundamental proposition must not, however, be misunderstood as though the
meaning were that we limited human beings, by seeing things in God, can now also
see through God's essence\footnote{S.~200. Mais il faut bien remarquer, qu'on ne
peut pas conclure que les esprits voyent l'essence de Dieu, de ce qu'ils voyent
toutes choses en Dieu de cette maniére. Parce que ce qu'ils voyent est
très-imparfait, et que Dieu est très-parfait.}. Far from it! How should the
imperfect be able to grasp the perfect?

``One more reason for assuming that we see all things because God wills it is also
this, that we thereby recognize our infinite dependence on God; for of ourselves
we know nothing. Non sumus sufficientes cogitare aliquid a nobis, tanquam ex
nobis, sed sufficientia nostra a Deo est. 2 Cor. 3, 5.''

``It is God himself who enlightens the philosophers in those cognitions which
these ungrateful people call natural, although they come to them from heaven. Deus
enim illis manifestavit. Rom. 1, 19. It is he who is the mind's true light: Pater
luminum. Jac. 1, 17. It is he who teaches men science: Qui docet hominem
scientiam. Psalm. 94, 10. In a word: it is the true light which enlightens all
those who come into this world: lux vera, quæ illuminat omnem hominem venientem in
hunc mundum. Joh. 1, 9.''

No one will be able to mistake that this is all theological, truly theological.

These features of a doctrine of knowledge untenable in itself can, however, in
this connection interest us only insofar as they furnish the presuppositions for a
corresponding view of nature, and thereby make evident what
% p. 17
lies in theology's concept of nature, when the consequences contained in it are
set forth as clearly, as many-sidedly, as unreservedly, as in Malebranche.

If God's revelation, continuing itself in the thinking minds, is the sole true
source of knowledge, then it presumably follows of itself that the essence of
things depends solely on the divine will, that nothing in nature has force and
power to act of itself, but that it must be God, God alone, who works all in all
things.

``Aristotle\footnote{S.~383. Car les termes ordinaires à ce Philosophe ne peuvent
servir qu'à exprimer confusément aux sens et à l'imagination les sentimens confus
que l'on a des choses sensibles: ou à faire parler d'une maniére si vague et si
indéterminée, que l'on n' exprime rien de distinct. Presque tous ses ouvrages,
mais principalement ses huit Livres de Physique \dots\ ne sont qu, une pure
Logique.} and the philosophers who follow his banner miss the truth and walk in
darkness. Relying on sense-impressions, they give us only empty words and confused
representations of things.''

``All the categories (termes) that designate only sensuous representations are
ambiguous --- ambiguous, be it noted, through error and ignorance, and thus in
turn the cause of an infinite multitude of errors.''

``The word `Ram,' for example, is ambiguous; it means a ruminant animal; it also
means a constellation which the sun enters in spring. By such an expression, to be
sure, no one is so easily deceived; for one would have to be an astrologer in the
extreme to imagine that there could here be any real connection between Ram and
Ram. But as regards the categories of sensuousness otherwise, one is easily
deceived; for there is almost no one who notices that they are one and all
ambiguous.''

``When, for example, the philosophers say: fire is hot; the plant is green, sugar
sweet, and so on, they understand it like children and everyday people, as though
it were the fire that contained the heat they feel, the plant the colour they
think they see, and the sugar the sweetness they taste when they eat it; and
% p. 18
so with all things we see or sense. They do not say two words without saying
something false: all that they say about this matter is obscure and confused.''

This is then proved on various grounds:

``First, human beings are highly at variance about the sensuous impressions: what
one finds sweet, another finds bitter; what one finds warm, another finds cold.''

``Second, the most diverse objects can make the same sense-impression. Gypsum,
snow, sugar, salt, and so on give the same sensation of colour: nevertheless the
whiteness of these things is very different, if one will otherwise judge by
something other than mere sensation. When one therefore says: the flour is white,
one says nothing definite.''

``Third, the corporeal properties that cause the most different sense-impressions
are now and then almost the same, and, conversely, those that produce roughly the
same impression are often very different. The qualities sweetness and bitterness
are almost not different, but the sensations of sweetness and bitterness are
essentially different. The motions (mouvemens) that cause pain and tickling differ
only by a more and a less, but the feelings of tickling and pain are essentially
different.''

``Another series of ambiguous expressions is drawn from logic, by whose general
concepts one can easily explain all things without knowledge of the things
themselves.\footnote{S.~386. Aristote est celui qui en a le plus fait usage, tous
ses Livres en sont pleins, et il y en a quelques uns, qui ne sont que pure
Logique. Il propose et résout toutes choses par ces beaux mots de \emph{genre},
d'\emph{espece}, d'\emph{acte}, de \emph{puissance}, de \emph{nature}, de
\emph{forme}, de \emph{facultez}, de \emph{qualitez}, de \emph{cause par soi}, de
\emph{cause par accident}. Ses sectateurs ont de la peine à comprendre que ces
mots ne signifient rien \dots}''

``One becomes no wiser by hearing the Aristotelians say that fire melts metals
because it has the faculty of melting, or that someone cannot digest because his
stomach has not the faculty, and so on.\footnote{S.~386. Le feu échauffe, séche,
durcit, et amollit, parce qu'il a la faculté de produire ces effets. Le sené purge
par sa qualité purgative, le pain même nourrit, si on le veut, par sa qualité
nutritive \dots\ Telles ou semblabes maniéres de parler ne sont point fausses,
mais c'est qu'en effet elles ne signifient rien.}''
% p. 19
``But what is the most dangerous thing about these philosophical turns of phrase
--- quality, faculty, substantiality, and the like --- is nevertheless this, that
one is thereby ensnared in the imagining that things themselves should contain the
true and original causes of the effects we see before our eyes.''

``It is a sinful error; for if there are in things themselves real causes, real
faculties to act of themselves, then there is in things something divine, and we
come then, just as the heathens, to deify nature.\footnote{S.~387. Que si l'on
vient ensuitte à considerer attentivement l'idée que l'on a de cause ou de
puissance d'agir, on ne peut douter que cette idée ne represente quelque chose de
divin. Car l'idée d'une puissance souveraine est l'idée de la souveraine divinité,
et l'idée d'une puissance subalterne est l'idée d'une divinité inférieure, mais
d'une véritable divinité, au moins, selon la pensée des Payens, supposé que ce
soit l'idée d'une puissance ou d'une cause véritable. On admet donc quelque chose
de divin dans tous les corps qui nous environnent, lorsqu'on admet des formes des
facultez, des qualitez, des vertus, et des êtres réels capables de produire
certains effets par la force de leur nature, et l'on entre ainsi insensiblement
dans le sentiment des Payens par le respect que l'on a pour leur Philosophie. Il
est vrai que la foi nous redresse, mais peut-être que l'on peut dire, que si le
coeur est chrétien, le fond de l'esprit est Payen.}''

``It is difficult to believe that one ought neither to fear nor to love the real
powers, the beings that can act upon us, that can punish us with pain and reward
us with joy. And since true worship consists in love and fear, it is furthermore
difficult to believe that one ought not to worship them.''

``Everything that can act upon us as a true and real cause is necessarily exalted
above us (est nécessairement au dessus de nous). When, therefore, God in Holy
Scripture proves to the Israelites that they ought to worship him --- that is, that
they ought to fear and love him --- the most important grounds he adduces are drawn
from his power to reward them and punish them.
% p. 20
He wills that they should honour him, since it is he alone who is the true cause of
good and evil; and since, according to the Prophet's word, no evil happens in
their city without its being he (Jehovah) who does it.

``The natural causes are therefore by no means real causes of the evil they seem
to inflict upon us, since it is God alone who works in them, he alone whom we ought
to fear and love: Soli Deo honor et gloria.''

If this is not theological, truly theological, then what is theological?

``If we assume,'' it says further, ``the philosophers' false opinion (cette fausse
opinion des Philosophes), namely, that the bodies surrounding us are real causes of
the joys and evils we feel, then reason seems to justify a religion like that of
the heathens, and to sanction a universal corruption of morals (déreglement
universel des meurs).''

``Even if reason does not directly tell us that we should therefore straightway
worship onions and leeks, it nevertheless advises us to render them a certain
veneration in proportion to the benefit they can render us.\footnote{S.~388. Il est
vrai que la raison n'enseigne pas qu'il faille adorer les oignons et les porreaux
par exemple, comme la souveraine divinité, parce qu'ils ne peuvent nous rendre
entiérement heureux, lorsque nous en avons, ou entiérement malheureux lors que
nous n'en avons point \dots\ Mais (S.~389), si l'on ne doit pas rendre un honneur
souverain aux porreaux et aux oignons, on peut toujours leur rendre quelque
adoration particuliére: je veux dire, qu'on peut y penser et les aimer en quelque
maniére, s'il est vrai qu'ils puissent en quelque sorte nous rendre heureux: on
doit leur rendre honneur à proportion du bien qu'ils peuvent faire.}''

Inwardly outraged at the falsity in so wretched a philosophy (de la fausseté de
cette misérable Philosophie), Malebranche then develops the doctrine: \emph{that
there is only one true cause, since there is only one true God; that nature, or the
force of things, is nothing other than God's will; that natural things are not true
causes, but solely occasional causes (des causes occasionelles).}

% p. 21
Here, then, we have the so-called occasionalist system: systema causarum
occasionalium.

How this theory is really to be understood is made vivid by the following example:

``When one ball collides with another and sets it in motion, it nevertheless
communicates nothing; for it does not itself have the impression it communicates.
Meanwhile a ball does become a natural cause of the motion it communicates. A
natural cause is therefore no true and real cause, but rather an occasional cause,
which determines the Creator of nature to act in this or that way in this or that
case.\footnote{Jvnf. S.~390 med S.~575. Quand je vois une boule qui en choque une
autre, mes yeux me disent, ou semblent me dire, qu'elle est véritablement cause du
mouvement qu'elle lui imprime: car la véritable cause qui meut les corps ne paroît
pas à mes yeux. Mais quand j'interroge ma raison, je vois évidemment que les corps
ne pouvant se remüer eux-mêmes, et leur force mouvante n'étant que la volonté de
Dieu qui les conserve successivement en différens endroits, ils ne peuvent
communiquer une puissance qu'ils n'ont pas, et qu'ils ne pourroient pas mêmes
communiquer quand elle seroit en leur disposition. Car il est évident qu' il faut
une sagesse, et une sagesse infinie pour régler la communication des mouvemens
avec la justesse, la proportion, et l'uniformité que nous voyons.} The laws by
which motions are communicated are effective only insofar as God's will works in
them. God has created the world because he willed it: dixit et facta sunt. There
are therefore in the sensuous world no forces, no powers, no true causes, and one
ought not, by ascribing to bodies forms, faculties, and real qualities that they
do not have, to deceive oneself and offend the divine majesty.''

``God does not need instruments in order to act (Dieu n'a pas besoin d'instrumens
pour agir); it is only needed that he will it, in order that the thing may be. His
power is his will. But to communicate his power to a man or to an angel can mean
nothing other than to will that, when a man or an angel wills that this or that
body be set in motion, it then moves.''
% p. 22
God is therefore, according to the cited example, the true cause of the motion,
the angel's will only an occasional cause.

``To make this still more evident, let us suppose that God willed the opposite of
what some spirits, for example the demons, will; then one could nevertheless in
this case not say that God communicates to them his power, for they can accomplish
nothing of what they wish. Meanwhile these spirits' wills would nonetheless be
natural causes of effects they produced. Such bodies would then move to the right
only because these spirits wished that they should move to the left: they thus
occasion God to do the opposite of what they will, and their wills thus also
become occasional causes.\footnote{S.~392. Il n'y a donc qu'un seul vrai Dieu et
qu'une seule cause qui soit véritablement cause: et l'on ne doit pas s'imaginer que
ce qui précéde un effet en soit la véritable cause. Dieu ne peut même communiquer
sa puissance aux créatures, si nous suivons les lumiéres de la raison: il n'en peut
faire de véritables causes: il n'en peut faire des dieux. Mais quand il le pouroit,
nous ne pouvons concevoir pourquoi il le voudroit. Corps, esprits, pures
intelligences, tout cela ne peut rien. C'est celui, qui a fait les esprits, qui
les éclaire et qui les agite... Enfin c'est l'Auteur de nôtre être qui exécute nos
volontez, semel jussit, semel paret. Il remue même nôtre bras lorsque nous nous en
servons contre ses ordres, car il se plaint par ses Prophetes que nous le faisons
servir à nos desirs injustes et criminels.}''

\begin{center}---\end{center}

In the history of theology there is no dogmatician who has held fast to the concept
of God's omnipotence with greater rigour than Malebranche, none who has expressed
with greater scientific consistency what lies in the dogma of creation than
precisely Malebranche.

He has not, like a Calvin, confined himself to a logical form for a religious
language of power; as a philosophical thinker he has everywhere striven to make the
truth of faith accessible to reason, and for that reason has not shrunk from
engaging with the real phenomena, and giving us an interpretation, consistent with
his theory, of the sensuously present order of things.

To be sure, Malebranche has not been able in a speculative way to see
% p. 23
into the Father's primal ground, to discover world-germs in the Son's Pleroma, and
to mediate the process of creation in the plastic power of the Spirit; but the
dogma he has, the pure dogma of creation out of nothing. Concerning this dogma he
has not merely declared that it must be thought (for because it is now and then
said of a dogma that it must be thought, it is not thereby said that it really gets
thought); no, he has in earnest attempted to think it, to hold the thought fast,
and to apply it to a scientific understanding of all that is created.

To apply the thought of creation to a scientific understanding of the created is to
see the whole world and each single thing in the world under the twofold point of
view: \emph{nothing and yet really something!}

The concept in which this contradiction is resolved is the concept: omnipotence;
the created is in
% p. 24
itself a nothing, but at omnipotence's command it has become something.

He who now looks at what has come to be so one-sidedly that he sees only its
nothingness --- for him the world is semblance and shadow. Whether theology will
call him a pantheist or an acosmist is a matter of indifference; the point is that
he denies creation in denying omnipotence.

He who looks at what has come to be so one-sidedly that he sees only its reality,
forgets its nothingness, and ascribes to it substantiality --- for him
coming-to-be is a natural process, and the world a realm of nature. Whether
dogmatics must otherwise reckon him among deists, naturalists, or atheists is in
this connection a matter of indifference; the point is that he denies creation in
denying omnipotence.

He, on the other hand, who in all relations under all lights constantly sees what
has come to be from the twofold side: as nothing in itself and yet as really
something, he sees in all things God's omnipotence, he thinks, though only in a
human way, the thought of creation, he holds fast to the dogma of creation.

This is what Malebranche has done; and in the libraries of theology one will surely
have to search long before one finds, as regards this point, a thinker as honestly
consistent as he.

Created out of nothing, all existing things are for him in themselves nothing, and
continue
% p. 24
to be so; created by omnipotence, existing things are for him really present in a
real world, which the Creator conserves and governs.

What is matter? The extended, which seems to be something when seen sensuously;
but, seen with the eye of reason, matter is essenceless, without form, without
faculty, without power to act in any way whatsoever in any direction whatsoever.
Nevertheless one cannot, according to the theory, say that matter does not exist;
on the contrary! it really exists; for it is created: the earnestness of the
almighty will vouches for its existence.

What is the soul? The thinking thing, which in finite thoughts thinks on itself,
and imagines itself to be something. But the soul, the finite spirit, is in itself
a nothing, a pure nothing, just as inactive, dead, and powerless as matter.
Nonetheless the soul really exists; for it is created, and thus, at omnipotence's
command, has become something out of nothing.

As the coming-to-be of created things is a miraculous transition from nothing to
something, so too their existence is a miraculous hovering between something and
nothing. Corps, esprits, pures intelligences tout cela ne peut rien. C'est celui,
qui a fait les esprits, qui les éclaire et qui les agite.

What holds of things themselves holds then also of their mutual relation. This
relation cannot possibly be essential, since things themselves are without
substantiality. Here there is no interaction between existing things: the natural
causes are, after all, not true causes. Nonetheless there is here, in this real
world, a real coherence, real laws, a real order of things; even if created things
are not true causes, they are yet, for the creating will, existing occasional
causes: the divine laws are revelations of the highest reason.\footnote{S.~390.
Dieu a créé le monde... et produit ainsi tous les effets que nous voyons arriver,
parce qu'il a voulu aussi certaines loix, selon lesquelles les mouvemens se
communiquent à la rencontre des corps: et parce que ces loix sont efficaces, elles
agissent, et les corps ne peuvent agir.}
% p. 25
It has indeed now been supposed that one could go a step further than Malebranche,
and improve his doctrine of creation by ascribing to the created some independence
after all, some capacity for inner self-development, some faculty for activity; it
is then especially Leibniz who has here wished to correct, or rather to push back,
the ``miracle theory'' by inventing the doctrine of monads, devising harmonia
præstabilita, and asserting principium rationis sufficientis.

It is, however, as regards the dogma, a dubious improvement. The created is then,
in creation, supposed to have received not merely a real existence, but at the same
time an existence in relative independence.

It is now this relative independence that Malebranche could not conceive, and
neither, when we look more closely, could a Leibniz. For what is relative
independence? A degree of existence outside omnipotence, a self-acting which,
however slight it may otherwise be, is nevertheless not omnipotence's acting, and
thus: a real limit to omnipotence.

If this real limit is now meant in earnest, then omnipotence is a really limited
power, and no longer omnipotence. If, on the other hand, the Almighty still, as
before, has the limit within his power, then the real limit is still, as before, a
nothing in itself, and the miraculous relation between Creator and creature must
continue.

A relation that has begun with a miracle, and been founded by a miracle, must also
be continued by an unceasing miracle: not merely for a few moments, not merely, if
I may so put it, some 8 to 14 days after the creation; but for centuries, for
millennia, in the eternity of eternities, created things must continue to be in
themselves what they originally are and have been, a nothing, a pure nothing for
the Almighty.

It is this consequence that Malebranche holds fast, and therefore he cannot
conceive self-activity in any creature whatsoever, be it angel or man, for
essential self-activity is a divine matter, and God himself cannot create gods: il
n'en peut faire de véritables causes, il n'en peut faire des Dieux.
% p. 26
The improvement that Leibniz introduces consists now in this, that he lets God
create gods\footnote{Chaque esprit étant comme une petite divinité dans son
département. Monadol. --- Les esprits étants comme de petits Dieux. Syst. nouv.};
for his relative monads are indeed like a swarm of little gods. Each of these
little gods carries, in the manner of gods, the possibility of a world within
itself, is, in the manner of gods, sufficient unto itself, leads, in the manner of
gods, an eternal life, and, even if it be weakened at a given time, can
nevertheless in all eternity not (unless it were by an extraordinary miracle) lose
the faculty of self-development\footnote{Atque ideo etiam nulla datur generatio, nec
mors perfecta rigorose loquendo. Sunt emim evolutiones et accretiones, quas
generationes appellamus, quemadmodum involutiones et diminutiones, quod mortem
vocamus. Princip. Philos. Thesis 76.}.

But what is now gained by putting all this self-activity into the created? Two
half-thoughts instead of one whole: consistency is broken. On the one hand,
creation is still supposed to be meant in earnest --- that is, the thought that
omnipotence's limit is in itself a nothing. On the other hand, the created's
independence is at the same time supposed to be meant in earnest, along with the
thought that omnipotence's limit is in itself a something. Omnipotence is thus
supposed at one and the same time to be both almighty and yet not almighty; behold,
that is what has been gained.

If this, what Leibniz would introduce, even when given a more comprehensible and
vulgar form, really is a progress; then indeed it cannot be denied that theology
necessarily must have something to say in scientific matters; for now the relation
between God and the world becomes as entangled as possible.

Now omnipotence has a limit, and no miracle happens; now omnipotence has no limit,
for a miracle has indeed happened. Here there is much to consider, when it --- be
it noted --- is to be considered scientifically. Here is nihil negativum, and here
is nihil privativum; here is omnipotentia, here is præscientia, here is providentia,
here is conservatio; here is conservatio rerum simplicium, here is annihilatio
rerum simplicium, here is conservatio nexus
% p. 27
cosmici, here is destructio nexus cosmici; and then here is besides concursus, and
concursus again both universalis and specialis and specialissimus.

If there then come to this a multitude of scriptural passages, well exegeted, well
cited, well grouped scriptural passages; if these passages are then, in a suitable
scheme, accompanied by critical-dogmatic remarks; if these remarks then also give
hints that those supernaturalist forms probably do not signify much; if these hints
are then so plain that everyone must be able to see where the limit is, and where
the limit is not, and consequently to discern for himself both what must be thought
and what shall be thought, and that here really nothing at all is thought: then we
have a tolerably fitting impression, proportionate to the importance of the matter,
of the treatment which the concept of nature, smuggled in by ``the creature's
self-activity,'' undergoes in theology.

By wishing to improve the ``miracle theory,'' Leibniz loses his scientific bearing,
and places himself in the ranks of those who otherwise quack about with concursus
and nexus cosmicus and destructio nexus cosmici\footnote{Omnis autem monas est
inexstingvibilis, neque enim substantiæ simplices nisi \emph{creando vel
annihilando}, id est miraculose oriri aut desinere possunt. Epist. ad Bierlingium.}
and the like; only that he does it with talent and diplomatic suppleness.

For what does it really mean that for everyday use no miracles happen, but only on
extraordinary occasions? It means, plainly understood, in plain words: for everyday
use God is not almighty; but on extraordinary occasions he is almighty.

That created things are in themselves nothing, and their natural coherence in
itself nothing --- this thought we human beings cannot endure for everyday use; it
comes into force only on isolated extraordinary occasions; in order now not to be
overwhelmed by the representation of God's omnipotence, omniscience, and
omni%
% p. 28
presence, one must push back the all-too-obtrusive miracle theory, and patch up the
reality of things with a little commercium substantiarum, harmonia præstabilita, a
little nexus cosmicus, principium rationis, and other such mediacies.

``Conservation is a continued creation''; if this thought is really held fast, then
Malebranche is right. There can then be no possible difference between creation and
conservation other than that which lies in the two concepts: beginning and
continuation.

The same Almighty who of old said: Become! must still, in every second, in every
moment, keep on saying: Continue to be!

If one nevertheless will now have another and ``more thorough'' difference between
creation and conservation, where do we then get to? Then we arrive at the result
that creation is the immediate, conservation the mediate miracle. But in this
``thoroughness'' there is no sense; for a mediate miracle is a contradiction. Dieu
n'a pas besoin d'instrumens pour agir, il suffit qu'il veüille.

Theology is too inclined to confuse the categories: one-sidedness and consistency.
When the thinker now and then seizes one of its own presuppositions, and develops
what is essentially contained in it, then at once it is said: what one-sidedness!

Such dismissals are unscientific. In science one whole thought, which really gets
thought, is worth more than ten half-thoughts, which on account of their inner
incompatibility never get thought, however much trouble one takes to explain that
they ``must be thought.''

Malebranche is, with respect to the dogma of creation, not a one-sided but a
consistent thinker. Creation is to him in earnest; conservation is therefore also
to him in earnest: he does not think pro forma! As little as what has come to be
could of old create itself, just as little can it even to this day conserve itself;
as little as what has come to be could originally be
% p. 29
created by another thing that has come to be, just as little can even now the one
thing that has come to be be either conserved or moved by the other\footnote{S.~391.
Je dis de plus qu'il n'est pas concevable, que Dieu puisse communiquer aux hommes ou
aux Anges la puissance qu'il a de remüer les corps; et que ceux qui prétendent, que
le pouvoir que nous avons de remüer nos bras, est une véritable puissance, doivent
avoüer que Dieu peut aussi donner aux esprits la puissance de créer, d'anéantir, de
faire toutes les choses possibles, en un mot, qu'il peut les rendre tout-puissans.}.

According to Malebranche, the dogma of creation cannot possibly rest on the
contradiction that God, in the beginning of time, should have used his omnipotence
to set himself an essential barrier, a barrier that is in itself a something, a
barrier that he did not have at every moment absolutely and unconditionally within
his power; no, creation is known by this, that God's omnipotence, out of love and
with freedom, sets itself an opposite, and through this opposite, at all points, in
all directions, at all times, relates itself to itself, to itself alone.

Therefore there is here a real world, and this world's existence is God's work;
therefore there is here in the world a real order and effective laws, for this
order is God's wisdom, these laws are God's thoughts. Therefore the law and the
corresponding occasional causes are meant in earnest --- so much in earnest that
Malebranche sees God's freedom precisely in this, that the Almighty, who at every
moment works every coming-to-be and every motion, binds himself by his own
thoughts, his own will, his own laws. Semel jussit, semel paret. Il remue même
notre bras, lorsque nous nous en servons contre ses ordres, car il se plaint par
ses Prophetes que nous le faisons servir à nos desirs injustes et criminels.

The truth in the natural is not nature; the natural is not natural: it is a work of
God.

The occasionalist theory carried through can be regarded as a scientific exposition
of the Psalmist's words (Ps. 104, 24): \emph{Lord, how manifold are thy works, thou
madest them all in wisdom.}
% p. 30
We have therefore not done theology an injustice in presenting Malebranche as an
example of theological scientificity.

\begin{center}---\end{center}

The \emph{doctrine of creation} carried through in theology is a carried-through
\emph{negation of nature}.

The theological concept of nature is negative; if this negative concept has
scientific validity, then natural science is not merely a very subordinate but also
a very dangerous science, for natural science's concept of nature is positive.

What is comprehensible, what is immediately evident, what is unmistakable, what is
incontrovertible --- of all such we say: ``it is natural.'' That a stone thrown into
the air must fall to the earth is natural; that he who holds his finger in the flame
must get it burned is natural; that he who falls into the water must get wet, that
he who cannot draw breath is suffocated, and so on, is all natural. It must be so;
it cannot possibly be otherwise. Why not? Let one give ever so many explanations and
adduce ever so many proofs: the nerve in all proofs, the explanatory element in all
explanations, is nevertheless the one thing: ``yes, it is natural!''

But in all that is natural nothing can be more natural than nature; one must
therefore not be surprised that mankind, despite the ``higher light'' with which
theological scientificity has already for centuries shone for the generations, still
to this day is so involuntarily caught in the imagining that nature really is, that
there is in existence something which, in the natural sense, we must call the
essence of nature.

The essence of nature! It would perhaps fall just as naturally to say: nature and
essence; for these two concepts are so intimately related that the one must at
bottom explain the other. Every grounding in science aims at demonstrating what must
necessarily follow from ``the essence of the matter,'' but what follows from ``the
essence of the matter'' is only that which lies in ``the nature of the matter.''

% p. 31
If now both concepts are at bottom one, and if ground and grounding are decisive
for all science; then one cannot be surprised that the human mind still constantly
holds fast to the thought that there nevertheless is and must be truth in a science
whose whole aim is to know nature in nature's light, although natural science's
natural grounds have doubtless come into a highly precarious conflict with
theology's ``deeper grounds.''

If natural science is to count as true science, then one must begin with the
presupposition \emph{that there is a natural reciprocal relation between the senses
and the sensuous world.}

But if this presupposition is to hold, theology must see to finding itself a place
outside science; for its scientificity, as theology, rests precisely on a negation
of this presupposition.

He who would know nature must rely on the testimony of the senses. To rely on the
testimony of the senses is not the same as to judge objects by the first, casual,
immediate sense-impression. To what extent natural science, in this respect, knows
how to make a distinction between semblance and essence, it has surely shown
sufficiently, among other things, in its defence of the Copernican system.

To rely on the testimony of the senses means to rely on the inner, essential
agreement between the sensuous and that which is sensed, that is: on the natural
truth of the sense-impression.

``It is a truth that the senses sense,'' with this one must begin in natural
science; ``it is an untruth that the senses sense,'' with this one must begin in
theology.

It is, scientifically understood, not a one-sidedness against a one-sidedness that
we have here before us; it is principle against principle, science against science,
that is at stake.

By relying on the senses a human being can, even with the best will, never bring it
further than to seeing nature in a natural light.
% p. 32
The more the senses count, the less faith counts; the less the senses count, the
more faith counts.

There is an infinite difference between the trust in the testimony of the senses
and of experiences that is demanded in natural science, and the trust in the
scientific assistance of the Holy Spirit that the theologians call faith.

If one nevertheless will extend the concept of faith to embrace every assumption
resting on inward reliance and trust, one still cannot deny that here there is a
difference between faith and faith. The natural scientist's faith is affirmation of
nature a priori; the theologian's faith is negation of nature a priori.

In science (it is not forgotten here that it is in science) either the natural
faith must annihilate the supernatural, or vice versa.

To wish to see nature at once \emph{both} in a natural and yet supernatural light,
\emph{both} in a holy and yet scientific light --- this willing-of-both is no sign
of speculative profundity, but a sign of thoughtlessness; for it is a contradiction.

If the natural light is to shine, the supernatural must be extinguished; if the
holy light is to rise, the scientific must be extinguished.

C'est pour nous accommoder à la maniére ordinaire de parler, que nous dirons dans
la suite, que les sens sentent; with these words Malebranche extinguishes the
natural light; the pious thinker ``beholds all things in God'': fides præcedit
intellectum.

If theology now really has such trust in its fides præcedit intellectum that with
it, in our days, it dares in earnest to defy natural science; then we shall not
merely concede that it has a faith, but even willingly grant that its faith goes
far in advance of its understanding.

If natural science is to be regarded as true science, then the principle must be
held fast \emph{that the natural understanding is in truth able to know the laws of
nature.}

But if this principle is to hold in science, theology
% p. 33
doubtless does best to seek itself a place outside science; for it is a principle
that faith must deny.

If faith is to judge, then the laws of existence are God's thoughts, his will the
only true cause. Il n'y a donc qu'un seul vrai Dieu et qu'une seule cause, qui soit
véritablement cause.

If the understanding is to judge, then the laws of existence are grounded in the
essence of things, and there must consequently be just as many true causes as there
are natural forces.

If faith is to judge, the laws of things must be capable of alteration at whatever
moment it pleases God; if the understanding is to judge, even the smallest
deviation from the law must be assumed to conflict with ``the nature of the
matter.''

Here judgment goes against judgment, and the essence of nature is at stake. When
the one judgment is heard, the other cannot be heard; one cannot at one time hear
two judgments because one has two ears.

The ear of faith cannot at all hear the judgment of the understanding, as little as
the ear of the understanding can hear the judgment of faith. The judgment of faith
is heard in the ``Oratory,'' but the judgment of the understanding must be heard in
the ``Laboratory.'' To place oneself midway between the ``Oratory'' and the
``Laboratory,'' prick up both ears, and say: credo, ut intelligam, in order to
listen to both sides, is to give oneself une fausse position.

If the judgment of the understanding is to be heard in science, the judgment of
faith must be heard outside science; one cannot reconcile such contradictory
judgments in one scientific judgment.

In order to conceal the glaring contradiction, one has attempted a humane
accommodation.

``The laws of nature are certainly God's thoughts; but God cannot, after all, alter
these thoughts, since it conflicts with his immutability.''

So runs the humane accommodation; but the truth presses on: the heterogeneous will
be separated.

It goes with humane accommodation for a time, as with Jesuitical
% p. 34
torture for a time; suddenly comes --- the moment! Spirits grow wroth, principles
are incensed: Luther is coarse, Galileo stamps on the floor.

If natural science is to be regarded as a science, then the idea of reason holds
\emph{that the realm of nature is a whole closed within itself, whose heterogeneous
fundamental parts, according to their qualities, are combined and separated in such
a way that the whole's particular members are formed, shaped, and dissolved with
inner necessity.}

But if this idea of reason is to hold in science, theology does best to move out of
science; for it lies precisely in its supernaturalist character to deny this idea.

``Ground,'' ``primal ground,'' ``becoming,'' ``birth,'' ``inner forming,''
``involuntary expression of one's own disposition,''\footnote{``Unsere Erklärung
von dem Natürlichen,'' says Link (Propyläen der Naturgeschichte. Berlin 1839, S.~4)
``ist dieselbe, welche Aristotoles giebt,'' and cites on this occasion Physic. lib.
2, c. 1.

Τὰ μὲν γὰρ φύσει ὄντα πάντα φαίνεται ἔχοντα ἐν ἑαυτοῖς ἀρχὴν κινήσεος καὶ στάσεος,
τὰ μὲν κατὰ τόπον, τὰ δὲ κατ' αὔξησιν καὶ φθίσιν, τὰ δὲ κατ' ἀλλοίωσιν.} ``instinct''
and ``impulse,'' are the categories that hover before consciousness when the talk
is of maternal nature.

That there is a truth in this, theology too acknowledged in its own way, when it
really began to speculate over the paternal nature.

In the doctrine of the Son's eternal birth from the Father, Athanasius was after
all proved right against Arius, and why? Because reason, on this point, was allowed
to discern that, if the Son was in unity of essence (ὁμοούσιος) with the Father,
then he could not have been created with freedom, but must have been born of
necessity.

But, in that theology thereby gained an insight into the essence of nature, its
scientific progress consisted really in this, that it transferred categories from
the natural to the supernatural. In its striving to comprehend the supernatural
nature (θεία φύσις) it thus barred its own way to an understanding of the natural
nature.
% p. 35
The Son's nature is born, but not created; Adam's nature is created, but not born;
Adam's children, on the other hand, are both created and born: can anything
scientific come out of this?

God cannot create gods, therefore the Son is born from eternity; that is what
Athanasius has discerned.

God cannot create gods, therefore there is no nature in the created; that is what
Malebranche has discerned.

What is in truth to become a nature must come to be in a natural way; what has come
to be by a miracle must have a miraculous existence, and can never become anything
in truth natural, just as little as it can ever become a natural true cause of
anything else natural. ``In the natural relations of nature there is just as little
truth as in natural nature''; this is the theological-scientific consequence of the
dogma of creation.

If Adam is created, then Adam is no nature, then there is no nature in Adam. If
there is no nature in Adam, then there is likewise no nature in Eve. If there is no
nature in Adam and Eve, then Adam and Eve's children likewise cannot be born; they
must all be created. In this sense one can then say with truth that every human
being is a creature of God.

Human nature, human birth, all such is, according to Malebranche, just as surely
semblance and sense-deception as it is semblance and sense-deception that the one
body should really be able to set the other in motion.

That a human being can be born can be thought; it is a pure, natural matter, and
there is to that extent unity in the thought. That a human being can be created can
be thought; it is a pure, supernatural matter, and there is to that extent unity in
the thought. But that a human being should be able \emph{both} to be born and to be
created, \emph{both} to be created and to be born --- this both-and cannot be
thought, for it is a contradiction.

If ``creationism'' and ``traducianism'' really let themselves be mediated by credo,
ut intelligam, then circle and conscience probably also let themselves be
% p. 36
mediated by credo, ut intelligam. A ``birth-process'' that is at the same time ``a
creation-process,'' a birth-creation, a creation-birth, a traducian creationism, a
creationist traducianism, a natural miracle-process, a miraculous nature-process:
if this is scientific, then natural science is unscientific.

Everyone who knows what science is must really admit that the theology of our day
finds itself in a peculiar position. In the same degree that it approaches the
supernatural, it must distance itself from nature; in the same degree that it will
approach a knowledge of nature, it must deny the supernatural.

Admonished, especially by Kant, the theologians for a time proceeded according to
the principle that everything depended on holding to the law of nature and to
practical reason.

``If only we could, in a scientific way, get rid of the many unreasonable
miracles!'' thought the Rationalists, and so they began to explain away the many
``unreasonable'' miracles. It succeeded, even beyond expectation; one miracle
disappeared after another --- in a scientific way. At last only one miracle
remained, namely creation; this miracle one might indeed reinterpret in particulars,
but on the whole not explain away; for the miracle of creation was after all at
bottom a quite reasonable miracle. By this reasonable procedure theology became
undeniably more and more reasonable, more and more scientific; at last only one
thing was lacking: consistency.

In science a consequence may sometimes keep one waiting, but never fail to come. The
Straussian critique was Rationalism's consequence; the consequence took away the
last remnant of the last miracle, and theology went backward out of science.

Chastened by the scourge of criticism, theology has again, little by little, begun
with credo, ut intelligam, to turn toward the supernatural; we cannot but take note
of its conversion. But a true conversion also has its consequence: ``show me thy
faith by thy works!''

If theology really makes it its earnest to turn toward the super%
% p. 37
natural, and carries through the consequence of conversion; then it shall never more
dread the miracle, never more bow to a law of nature.

If theology really makes it its earnest to turn toward the supernatural, and carries
through the consequence of conversion; then its powers will grow, then it will
become strong: strong in Credo, strong in Intelligam, strong in the miracle! so
strong in the miracle that its miracle-strength strikes ``all-nature'' dead, and
goes victorious --- out of science.

If theology really makes it its earnest to turn toward the supernatural, and carries
through the consequence of conversion, then it obtains peace from all criticism;
then it is gathered to its fathers, and lands in a better harbour: Athanasius has
long since prepared it a place; Malebranche stands waiting in the ``Oratory.''

\begin{center}---\end{center}

Between theology and natural science there is no understanding: theology's concept
of nature is not natural, and natural science's is not theological.

\begin{center}---\end{center}

% ---- Appendix: Indbydelsesskrift (printed pp. 38--47) — OUT OF SCOPE ----
% Per Hans's decision (2026-07-04), the deliverable is the treatise (pp. 3--37) only.
% The appendix (the three doctoral vitae — Salomonsen, Müller, Helweg — and the
% Reformation-festival invitation ending "Under Universitetets Segl.") is NOT
% translated. The Danish remains available in transcription.tex, pp. 38--47.

\end{document}
